\label{sec:discussao}

\chapter{Discussão}

Esse capítulo apresenta uma análise crítica dos resultados obtidos, destacando as limitações técnicas e operacionais da solução \acs{open_source} desenvolvida, bem como as lições aprendidas durante o processo de implementação.

\section{Estimativa de Custos Operacionais}

A adoção de ferramentas de código aberto elimina os custos de licenciamento de
\textit{software}, que em soluções proprietárias equivalentes podem representar a parcela
mais expressiva do orçamento de um \textit{pipeline} de dados. No entanto, a ausência de
licenças não implica ausência de custos. Os gastos com infraestrutura, servidores,
armazenamento e transferência de dados persistem e devem ser estimados com precisão
para que o benefício econômico do \acs{open_source} seja quantificado de forma correta e justa.

\subsection{Premissas da Estimativa}

A comparação entre a solução de código aberto e os equivalentes gerenciados nas
principais nuvens públicas exige um conjunto de premissas explícitas para garantir
que o confronto seja reproduzível e justo.

\textbf{Região geográfica equivalente:} As três provedoras avaliadas,
Amazon Web Services (AWS), Microsoft Azure e Google Cloud Platform (GCP),
oferecem infraestrutura na Região Metropolitana de São Paulo. Para este trabalho,
foram selecionadas as seguintes regiões: sa-east-1 (AWS),
Brazil South (Azure) e southamerica-east1 (GCP), garantindo
paridade de latência e conformidade com requisitos de residência de dados no
Brasil \cite[]{aws2026pricing, azure2026pricing, gcp2026pricing}.

\textbf{Configuração mínima viável:} A execução simultânea dos contêineres Docker
que compõem o \textit{pipeline}, com Apache Airflow (agendador, servidor web e
executor), PostgreSQL, Apache Superset e os \textit{scrapers} Selenium, requer
no mínimo 2 vCPUs e 8 GB de RAM. Essa configuração é adotada como referência para
todas as instâncias de máquina virtual comparadas. 

\textbf{Volume de dados:} Cada execução completa do \textit{pipeline} coletou
3.084 KB de dados brutos em média. Projetando execuções diárias ao longo de 30 dias,
o volume mensal estimado é de aproximadamente 90 MB, incluindo dados transformados,
metadados do Airflow e histórico de execuções no PostgreSQL.

\textbf{Armazenamento em disco:} Será adotado 100 GB de disco SSD para acomodar
os dados coletados, o banco PostgreSQL com histórico de 12 meses de execuções,
os arquivos de log do Airflow e as imagens dos contêineres Docker.

\textbf{Janela de operação:} As instâncias de \textit{VM} são tratadas como instâncias
\textit{on-demand} ativas 24 horas por dia, 7 dias por semana (730 horas por mês),
sem planos de reserva (\textit{Reserved Instances}, \textit{Savings Plans} ou
\textit{Committed Use Discounts}).

Os valores de referência são aqueles publicados nas páginas oficiais de preços das três provedoras, na modalidade \textit{pay-as-you-go},
vigentes em maio de 2026, ano de publicação deste trabalho \cite[]{aws2026pricing, azure2026pricing, gcp2026pricing}.
Todos os valores são expressos em dólares norte-americanos (US\$).

\subsection{Infraestrutura de Servidor}

A Tabela~\ref{tab:custo_vm} compara o custo mensal de uma instância de computação
com 2 vCPUs e 8 GB de RAM nas três provedoras, nas regiões sul-americanas de
referência.

\begin{table}[H]
    \centering
    \small
    \setlength{\extrarowheight}{3pt}
    \renewcommand{\arraystretch}{1.3}
    \begin{tabular}{
        >{\raggedright\arraybackslash}p{2.6cm}
        >{\centering\arraybackslash}p{1.1cm}
        >{\centering\arraybackslash}p{1.1cm}
        >{\raggedright\arraybackslash}p{1.8cm}
        >{\raggedright\arraybackslash}p{3.4cm}
        >{\centering\arraybackslash}p{1.6cm}
        >{\centering\arraybackslash}p{1.8cm}
    }
        \toprule
        \textbf{Instância} & \textbf{vCPU} & \textbf{RAM} & \textbf{Provedor} &
        \textbf{Região} & \textbf{US\$/hora} & \textbf{US\$/mês} \\
        \midrule
        \rowcolor[gray]{0.96}
        t3.large         & 2 & 8 GB & AWS   & sa-east-1            & 0,1344 & 98,11 \\
        D2s v3           & 2 & 8 GB & Azure & Brazil South         & 0,1590 & 116,07 \\
        \rowcolor[gray]{0.96}
        e2-standard-2    & 2 & 8 GB & GCP   & southamerica-east1   & 0,1064 & 77,67 \\
        \bottomrule
    \end{tabular}
    \caption{Custo mensal de instância de computação 2 vCPU / 8 GB nas regiões
    sul-americanas, modalidade \textit{on-demand}, Linux (maio/2026).}
    \label{tab:custo_vm}
\end{table}

A instância equivalente na GCP é aproximadamente 21\% mais barata do que na AWS e
33\% mais barata do que na Azure, para a mesma configuração de hardware, o que
reflete as distintas estratégias de precificação regional de cada provedora.

Para a alternativa de serviços gerenciados, o componente de maior impacto financeiro
é o orquestrador Apache Airflow gerenciado. A Tabela~\ref{tab:custo_airflow}
apresenta as tarifas de referência para esta categoria.

\begin{table}[H]
    \centering
    \small
    \setlength{\extrarowheight}{3pt}
    \renewcommand{\arraystretch}{1.3}
    \begin{tabular}{
        >{\raggedright\arraybackslash}p{5.0cm}
        >{\centering\arraybackslash}p{1.8cm}
        >{\raggedright\arraybackslash}p{3.6cm}
        >{\centering\arraybackslash}p{2.2cm}
    }
        \toprule
        \textbf{Serviço Gerenciado} & \textbf{Provedor} & \textbf{Região} &
        \textbf{US\$/mês} \\
        \midrule
        \rowcolor[gray]{0.96}
        Amazon MWAA\textsuperscript{(a)} (Small + 1 worker) &
        AWS & us-east-1\textsuperscript{(a)} & 397,85 \\
        Apache Airflow via Azure Container Apps (estimado) &
        Azure & Brazil South & aprox. 100,00 \\
        \rowcolor[gray]{0.96}
        Cloud Composer 2 (ambiente pequeno) &
        GCP & southamerica-east1 & aprox. 200,00 \\
        \bottomrule
    \end{tabular}
    \caption{Custo mensal do Apache Airflow gerenciado nas principais provedoras de
    nuvem (maio/2026).}
    \label{tab:custo_airflow}
\end{table}

{\footnotesize \textsuperscript{(a)} O Amazon MWAA (\textit{Managed Workflows for
Apache Airflow}) não está disponível na região sa-east-1 em maio de 2026.
o valor de referência é da região us-east-1, a menor tarifa disponível.
Inclui um ambiente \textit{small} (US\$ 0,49/hora) e um \textit{worker} adicional
(US\$ 0,055/hora), 730 horas/mês \cite[]{aws2026pricing}.}

\medskip

A diferença entre hospedar o Apache Airflow em uma \textit{VM} própria, incluso no custo
da instância, e contratar um serviço gerenciado equivalente é expressiva. Somente
no componente de orquestração, a solução gerenciada da AWS representa um gasto
superior a US\$ 4.700 por ano, mais de cinco vezes o custo anual da \textit{VM} inteira
na opção mais econômica (GCP: US\$ 77,67 por mês, o que equivale a US\$ 932 por ano).

\subsection{Armazenamento}

A Tabela~\ref{tab:custo_storage} apresenta as tarifas de armazenamento em disco bloco
(SSD) e em objeto para as três regiões de referência.

\begin{table}[H]
    \centering
    \small
    \setlength{\extrarowheight}{3pt}
    \renewcommand{\arraystretch}{1.3}
    \begin{tabular}{
        >{\raggedright\arraybackslash}p{2.6cm}
        >{\centering\arraybackslash}p{1.8cm}
        >{\raggedright\arraybackslash}p{3.8cm}
        >{\raggedright\arraybackslash}p{3.4cm}
        >{\centering\arraybackslash}p{1.8cm}
    }
        \toprule
        \textbf{Tipo} & \textbf{Provedor} & \textbf{Serviço} &
        \textbf{Região} & \textbf{US\$/GB/mês} \\
        \midrule
        \rowcolor[gray]{0.96}
        Bloco SSD & AWS   & EBS gp3                   & sa-east-1            & 0,112  \\
        Bloco SSD & Azure & Premium SSD v2 LRS         & Brazil South         & 0,152  \\
        \rowcolor[gray]{0.96}
        Bloco SSD & GCP   & Persistent Disk SSD        & southamerica-east1   & 0,170  \\
        \midrule
        Objeto    & AWS   & S3 Standard                & sa-east-1            & 0,0405 \\
        \rowcolor[gray]{0.96}
        Objeto    & Azure & Blob Storage Hot LRS        & Brazil South         & 0,0326 \\
        Objeto    & GCP   & Cloud Storage Standard      & southamerica-east1   & 0,035  \\
        \bottomrule
    \end{tabular}
    \caption{Tarifas de armazenamento por tipo e provedora nas regiões sul-americanas
    (maio/2026).}
    \label{tab:custo_storage}
\end{table}

Para os 100 GB de armazenamento em bloco (SSD) adotados como premissa, o custo
mensal resulta em US\$ 11,20 (AWS), US\$ 15,20 (Azure) e US\$ 17,00 (GCP).
O armazenamento em objeto, que é comumente utilizado nos serviços gerenciados para
\textit{data lakes} e ingestão de dados em escala, apresenta tarifas de 2,5
a 5 vezes menores que o armazenamento em bloco. Nessa categoria, o Azure Blob
Storage Hot LRS é o mais econômico entre as três provedoras avaliadas.

\subsection{Transferência de Dados}

A Tabela~\ref{tab:custo_transferencia} apresenta as tarifas de transferência de
dados de saída (\textit{egress}) para a internet nas três regiões de referência.

\begin{table}[H]
    \centering
    \small
    \setlength{\extrarowheight}{3pt}
    \renewcommand{\arraystretch}{1.3}
    \begin{tabular}{
        >{\centering\arraybackslash}p{1.8cm}
        >{\raggedright\arraybackslash}p{4.0cm}
        >{\centering\arraybackslash}p{2.8cm}
        >{\centering\arraybackslash}p{2.8cm}
    }
        \toprule
        \textbf{Provedor} & \textbf{Região} &
        \textbf{Franquia gratuita} & \textbf{US\$/GB (acima)} \\
        \midrule
        \rowcolor[gray]{0.96}
        AWS   & sa-east-1            & 1 GB/mês   & 0,250 \\
        Azure & Brazil South         & 100 GB/mês & 0,120 \\
        \rowcolor[gray]{0.96}
        GCP   & southamerica-east1   & ---         & 0,230 \\
        \bottomrule
    \end{tabular}
    \caption{Tarifas de transferência de dados de saída (\textit{egress}) para a
    internet nas regiões sul-americanas, modalidade \textit{on-demand} (maio/2026).}
    \label{tab:custo_transferencia}
\end{table}

Para o volume coletado pelo projeto deste trabalho, aproximadamente 90 MB mensais com
execuções diárias, o custo de transferência é negligenciável em todas as
provedoras. Contudo, em cenários com múltiplos usuários acessando \textit{dashboards}
ou com volumes de coleta substancialmente maiores, o \textit{egress} passa a ser
fator relevante. A tarifa da AWS (US\$ 0,250/GB) é 2,1 vezes superior à da Azure
(US\$ 0,120/GB), diferença que deve ser considerada em arquiteturas com consultas
frequentes a dados armazenados em nuvem.

\subsection{Custo Total e Comparação com Soluções Proprietárias}

A Tabela~\ref{tab:custo_total} consolida os custos mensais estimados para dois
cenários: (a) implantação da solução \acs{open_source} em uma \textit{VM} própria na
respectiva nuvem e (b) substituição de cada componente pelo serviço gerenciado
equivalente da mesma provedora.

\begin{table}[H]
    \centering
    \small
    \setlength{\extrarowheight}{3pt}
    \renewcommand{\arraystretch}{1.3}
    \begin{tabular}{
        >{\raggedright\arraybackslash}p{3.6cm}
        >{\centering\arraybackslash}p{2.2cm}
        >{\centering\arraybackslash}p{2.2cm}
        >{\centering\arraybackslash}p{2.2cm}
        >{\centering\arraybackslash}p{2.2cm}
    }
        \toprule
        \textbf{Componente / Cenário} &
        \textbf{AWS} \newline \footnotesize{sa-east-1} &
        \textbf{Azure} \newline \footnotesize{Brazil South} &
        \textbf{GCP} \newline \footnotesize{s.a.-east1} &
        \textbf{Referência} \\
        \midrule

        \multicolumn{5}{l}{\textit{Cenário A — Open Source auto-hospedado (US\$/mês)}} \\
        \midrule
        \rowcolor[gray]{0.96}
        VM (2 vCPU, 8 GB RAM)         & 98,11  & 116,07 & 77,67  & Tab.~\ref{tab:custo_vm} \\
        Armazenamento (100 GB SSD)     & 11,20  & 15,20  & 17,00  & Tab.~\ref{tab:custo_storage} \\
        \rowcolor[gray]{0.96}
        \textbf{Total — Cenário A}     & \textbf{109,31} & \textbf{131,27} & \textbf{94,67} & \\
        \midrule

        \multicolumn{5}{l}{\textit{Cenário B — Serviços gerenciados equivalentes (US\$/mês)}} \\
        \midrule
        \rowcolor[gray]{0.96}
        Orquestração (Airflow) & 397,85\textsuperscript{(a)} & aprox. 100,00 & aprox. 200,00 & Tab.~\ref{tab:custo_airflow} \\
        PostgreSQL gerenciado  & 110,78 & 104,23 & 121,68 & \\
        \rowcolor[gray]{0.96}
        Visualização (BI)      & 18,00\textsuperscript{(b)} & 13,70\textsuperscript{(c)} & 0,00\textsuperscript{(d)} & \\
        Transformações (dbt)   & aprox. 2,20\textsuperscript{(e)} & aprox. 5,00\textsuperscript{(f)} & 0,00\textsuperscript{(g)} & \\
        Armazenamento (100 GB) & 11,20  & 15,20  & 17,00  & Tab.~\ref{tab:custo_storage} \\
        \rowcolor[gray]{0.96}
        \textbf{Total — Cenário B}     & \textbf{aprox. 540} & \textbf{aprox.238} & \textbf{aprox. 339} & \\
        \midrule
        \textbf{Fator B / A}           & \textbf{4,9 vezes}  & \textbf{1,8 vezes}  & \textbf{3,6 vezes}  & \\
        \bottomrule
    \end{tabular}
    \caption{Comparativo de custo mensal entre a solução \acs{open_source}
    auto-hospedada (Cenário A) e os serviços gerenciados equivalentes (Cenário B)
    nas três principais provedoras de nuvem, nas regiões sul-americanas (maio/2026).}
    \label{tab:custo_total}
\end{table}

{\footnotesize
\textsuperscript{(a)} Amazon MWAA Small em us-east-1 (não disponível em sa-east-1):
ambiente US\$ 0,49/h + 1 worker US\$ 0,055/h por 730 h/mês.

\textsuperscript{(b)} Amazon QuickSight com um autor.

\textsuperscript{(c)} Microsoft Power BI Pro com um usuário.

\textsuperscript{(d)} Looker Studio no plano gratuito.

\textsuperscript{(e)} AWS Glue com 30 execuções mensais de 5 minutos (2 DPUs).

\textsuperscript{(f)} Azure Data Factory, com uso leve de Integration Runtime.
\textsuperscript{(g)} Dataform incluído sem custo no BigQuery.
}

\medskip

Os resultados evidenciam que a solução \acs{open_source} auto-hospedada custa
entre US\$ 94,67 (GCP) e US\$ 131,27 (Azure) por mês, enquanto os equivalentes
gerenciados variam de US\$ 238 (Azure) a US\$ 540 (AWS). O fator de custo, a
razão entre o Cenário B e o Cenário A na mesma provedora, oscila de 1,8 vezes (Azure)
a 4,9 vezes (AWS), chegando a 5,7 vezes quando se compara a opção \acs{open_source} mais
econômica (GCP: US\$ 94,67) ao equivalente gerenciado mais caro (AWS: US\$ 540).

A diferença menos expressiva no caso da Azure (1,8 vezes) decorre de um fator estrutural.
A plataforma não oferece um serviço de Apache Airflow totalmente gerenciado em
paridade com o Amazon MWAA ou o Cloud Composer do GCP. O equivalente Azure foi
estimado com base em Azure Container Apps, que é uma plataforma de contêineres
serverless, que reduz artificialmente o custo do Cenário B e não inclui
os custos operacionais de configuração, atualização e monitoramento que uma equipe
técnica precisaria absorver.

É necessário ressaltar que a comparação não é inteiramente semelhante. Os serviços
gerenciados oferecem acordos de nível de serviço (\textit{SLAs}) de disponibilidade,
escalabilidade automática, atualizações de segurança aplicadas pela própria
provedora e monitoramento integrado. Na solução \acs{open_source}, esses mesmos
fatores, como atualização de contêineres, reinicialização após falhas, \textit{backup}
do PostgreSQL, renovação de certificados TLS, são responsabilidade da equipe
interna. O Custo Total de Propriedade (\textit{Total Cost of Ownership}, TCO)
depende, portanto, do custo da hora de trabalho da equipe técnica. Para
organizações com equipes qualificadas, o \acs{open_source} permanece
substancialmente mais barato. Para organizações sem colaboradores técnicos dedicados,
o valor agregado dos serviços gerenciados pode justificar o custo adicional.

\section{Análise Crítica e Lições Aprendidas}

\subsection{Limitações de Escalabilidade das Ferramentas \acs{open_source}}

A solução desenvolvida demonstrou viabilidade técnica e econômica para o escopo
escolhido. Contudo, a escolha por ferramentas \acs{open_source}
auto-hospedadas implica limitações de escalabilidade que é necessário delimitar
com precisão, para que futuros usuários possam antecipar os pontos de quebra
da arquitetura.

O LocalExecutor do Apache Airflow opera exclusivamente no processo do agendador
e não distribui tarefas para múltiplos nós de computação. Para os 16 fluxos
sequenciais do protótipo essa limitação é irrelevante, mas qualquer ampliação
substancial no número de \textit{scrapers}, ou a necessidade de execução
paralela de múltiplas DAGs, exigirá migração para o CeleryExecutor ou o
KubernetesExecutor, ambos com custos de infraestrutura e configuração
muito maiores.

O PostgreSQL está implantado como contêiner único, sem replicação ou mecanismo
de \textit{failover} automático. Uma falha no volume de dados ou no próprio contêiner
interrompe imediatamente a disponibilidade do \textit{data warehouse}. Em cenários
de produção que demandem alta disponibilidade, seria necessário implementar
\textit{streaming replication} nativo do PostgreSQL ou soluções de orquestração de
\textit{failover}.

O Apache Superset, na configuração padrão com SQLite como banco de metadados,
não foi projetado para ambientes multiusuário de alta concorrência. Para
\textit{dashboards} acessados por muitos usuários simultâneos ou com consultas
SQL complexas sobre grandes volumes de dados, o tempo de resposta pode diminuir
progressivamente. A escalabilidade horizontal do Superset requer múltiplos
\textit{workers} Celery, um broker Redis dedicado e migração do banco de metadados
para PostgreSQL externo, o que aproxima a complexidade operacional da solução
\acs{open_source} com a de um serviço gerenciado na nuvem.

O Docker Compose foi projetado para orquestração em um único hospedeiro e não
suporta distribuição de contêineres entre múltiplos \textit{nodes} de forma nativa. A migração
para um ambiente multihost exigiria a inclusão de Kubernetes,
introduzindo uma nova camada de complexidade de infraestrutura e, consequentemente,
um aumento massivo no custo de manutenção.

Em conjunto, essas limitações não invalidam a abordagem para o escopo do projeto, 
dezenas de fontes, execuções diárias e volumes na casa dos gigabytes, mas
definem com clareza o horizonte da solução atual. O crescimento além
desses limites demandará ou a evolução das ferramentas \acs{open_source} para
configurações distribuídas, ainda que dentro do ecossistema de código aberto e com
custo substancialmente inferior ao dos serviços gerenciados, ou até mesmo a transição gradual para serviços gerenciados em
nuvem, cuja decisão deve ser orientada pelo líder técnico e pela capacidade da equipe
responsável pela manutenção.

\subsection{Síntese das Lições Aprendidas}

O desenvolvimento evidenciou que a viabilidade de um \textit{pipeline}
de dados \acs{open_source} para extração e análise de informações de mercado turístico
é técnica e economicamente comprovada. As principais lições aprendidas, com valor
imediato para implementações futuras, são organizadas abaixo.

As diferentes formas das fontes é o principal desafio de engenharia. Os 10
\textit{scrapers} desenvolvidos implementam ao menos seis estratégias distintas de
extração, \textit{scroll} infinito, carrossel, paginação, botão ``ver mais'',
\textit{dropdown} de navegação e navegação multi-nível, porque cada operadora turística
construiu seu site de forma independente. Essa diversidade torna impossível uma solução
genérica a baixo custo, e reforça a necessidade de manutenção contínua, alterações no \textit{layout}
de qualquer site exigem atualização do \textit{scraper} correspondente.

A deduplicação incremental é uma funcionalidade, não um detalhe. A estratégia
de comparar tuplas de campos-chave antes de inserir registros, implementada via
get\_existing\_records() no db\_manager.py, permite que os
\textit{scrapers} sejam reexecutados sem gerar duplicatas, o que é essencial tanto para
a confiabilidade dos dados quanto para a operação de recuperação após falhas.

O LocalExecutor do Airflow é suficiente para o escopo atual, mas cria débito
arquitetural. A escolha deliberada pelo LocalExecutor simplificou enormemente a
infraestrutura inicial, mas qualquer incremento substancial no número de
\textit{scrapers} exigirá migração para um executor distribuído. Documentar essa decisão
e seus \textit{trade-offs} desde o início é fundamental para que futuros mantenedores
compreendam o contexto da escolha, estejam cientes do débito técnico.

Ferramentas \acs{open_source} não eliminam custos, mas os redistribuem. O
custo de licenciamento zero é substituído por custos de operação (servidor), custo de
aprendizado (curva de adoção das ferramentas) e custo de manutenção (atualizações,
correções de \textit{scrapers}, monitoramento). Para organizações sem equipe técnica
dedicada, a comparação com soluções SaaS deve considerar o custo total de propriedade
(\textit{Total Cost of Ownership} ou TCO), e não apenas o custo de licença.

A separação dos dados em camadas demonstrou seu valor prático. O modelo de quatro camadas
implementado, \textit{scrapers}, passando pelo DuckDB/PostgreSQL \textit{raw},
depois pelos modelos dbt staging/marts, até o Superset, permitiu que cada componente
evoluísse de forma independente. A adição de novos \textit{scrapers} não afetou os
modelos dbt existentes, e a refatoração dos modelos não exigiu alterações nos
\textit{scrapers}. Essa independência é um resultado direto da adesão ao princípio de
separação de responsabilidades e ao padrão de dados em camadas (\textit{medallion
architecture} simplificada).

Em resumo, o \textit{pipeline} desenvolvido demonstra que é possível construir uma
infraestrutura analítica funcional, auditável e economicamente acessível inteiramente
com ferramentas de código aberto. As limitações
identificadas não invalidam a abordagem, mas delimitam com precisão seu envelope de
aplicabilidade. O protótipo é adequado para dezenas de fontes e volumes na casa dos
gigabytes, e requer evolução arquitetural planejada para os próximos patamares de escala.


